\documentclass[11pt]{article}
\usepackage[margin=1in]{geometry}
\usepackage{graphicx}
\usepackage{booktabs}
\usepackage{hyperref}
\usepackage{amsmath}
\usepackage{xcolor}
\usepackage{caption}
\usepackage{float}

\title{ReproBench-CXR: A Reproducible Benchmark Pipeline for\\
Chest Radiograph Nodule Detection on JSRT}

\author{
  Kamel Assaf \\
  \texttt{assaf.kamel@hotmail.com} \\
  Independent Researcher
}

\date{\today}

\begin{document}
\maketitle

\begin{abstract}
We present ReproBench-CXR, an open-source pipeline for reproducible
benchmarking of chest radiograph (CXR) nodule detection. Our contributions
are threefold: (1) a configurable classical baseline based on scale-normalized
Laplacian-of-Gaussian (LoG) blob detection with lung-field masking, contrast
and circularity gating, and rib-artifact suppression; (2) a fine-tuned Faster
R-CNN detector using a COCO-pretrained backbone, trained on the JSRT dataset
without any external CXR data; (3) a FROC evaluation harness with bootstrap
95\% confidence intervals and a web interface for visual inspection and
case sharing. On the JSRT dataset (247 images, 152 annotated nodules), our
fine-tuned Faster R-CNN achieves \textbf{0.708} sensitivity at 2
false positives per image, compared to \textbf{0.083} for the
classical LoG baseline, an improvement of \textbf{8.5}$\times$. At the
stricter operating point of 0.5 false positives per image, the classical
pipeline produces zero true detections, while the DL model achieves
\textbf{0.667} sensitivity. All code, model checkpoints, and the exact
train/val/test split are released to enable direct reproduction of every
number in this paper.
\end{abstract}

\section{Introduction}

Chest radiography is the most common medical imaging examination worldwide,
with over 2 billion studies performed annually. Pulmonary nodule detection on
CXR is a canonical computer-aided detection (CAD) task, and the Japanese
Society of Radiological Technology (JSRT) dataset has served as a standard
benchmark since 2000 \cite{shiraishi2000jsrt}. Despite this, the field suffers
from a reproducibility problem: published sensitivities at fixed false-positive
rates are difficult to compare because preprocessing, lung-mask definition,
and FROC computation differ across papers.

Recent work on chest radiograph analysis has focused on classification
\cite{habeeb2025modified}, segmentation \cite{thamilarasi2025enhanced},
and multi-resolution detection architectures \cite{li2020multi}. These are
valuable contributions, but their reported numbers are difficult to compare
across papers because of differing preprocessing, split definitions, and
metric implementations. Our contribution is orthogonal: a small, complete,
reproducible pipeline on a single dataset (JSRT) with all artifacts
released.

We present a small, complete, reproducible pipeline to address this. Our
contributions are:

\begin{itemize}
\item A configurable classical LoG-based baseline with explicit,
documented preprocessing choices (Section~\ref{sec:classical}).
\item A fine-tuned Faster R-CNN detector with COCO-pretrained weights but
no CXR-specific pretraining, demonstrating what transfer from natural-image
detectors alone can achieve on JSRT (Section~\ref{sec:dl}).
\item A FROC evaluation harness with bootstrap confidence intervals and a
fixed random seed for the train/val/test split (Section~\ref{sec:eval}).
\item A web interface for visual inspection and case sharing, so that
qualitative results can be verified independently
(Section~\ref{sec:tooling}).
\end{itemize}

We are explicit about what we do \textit{not} claim: our fine-tuned Faster
R-CNN is not state of the art on JSRT. Rather, our contribution is the
reproducibility of the pipeline and the honest reporting of confidence
intervals on a small test set.

\section{Related work}

\subsection{Classical nodule detection on JSRT}
LoG-based approaches dominated CXR CAD through the 2000s. van Ginneken et
al.~\cite{ginneken2001} proposed a multi-scale LoG framework that remains a
standard baseline, reporting sensitivities of approximately 60--70\% at 2--5
false positives per image. Our classical baseline (Section~\ref{sec:classical})
is a modern reimplementation with explicit gating and mask definitions, with
the goal of being directly comparable to those historical numbers.

\subsection{Deep learning on chest radiographs}

Recent deep learning work on chest radiographs has primarily focused on
image-level classification \cite{habeeb2025modified} or lung-field
segmentation \cite{thamilarasi2025enhanced}. These numbers correspond to
different tasks than ours. Habeeb et al.~\cite{habeeb2025modified} report
high classification accuracy on lung cancer detection, but classification
answers ``does this image contain a nodule?'' rather than ``where is the
nodule?''. Our task is localization, which we argue is the more clinically
relevant formulation for a computer-aided detection (CAD) tool.

The closest prior work to ours is Li et al.~\cite{li2020multi}, who propose
a multi-resolution CNN for lung nodule detection on chest radiographs.
Their approach is architecturally related to ours in that both use a
multi-scale feature representation, but Li et al. train from scratch on
their own dataset, whereas our contribution is a \textit{reproducible}
transfer-learning baseline on a public dataset, with explicit
confidence intervals on a fixed train/val/test split. Direct numerical
comparison is not possible because the datasets differ; the relevant
claim of our paper is reproducibility of the pipeline, not a
state-of-the-art number.

\section{Dataset and split}
\label{sec:data}

We use the JSRT dataset \cite{shiraishi2000jsrt}, which contains 247 chest
radiographs (154 with one or more nodules, 93 nodule-free). Nodule annotations
in the provided \texttt{CLNDAT\_EN.txt} give the pixel center and diameter of
each annotated nodule on a 2048$\times$2048 grid, along with a subjective
subtlety score from 1 (very subtle) to 5 (obvious).

\textbf{Split.} We split the 154 annotated images into a 70/15/15
train/val/test split (\textbf{107/24/23} images) with a fixed seed of 0.
Because there is only one annotation per image in the JSRT file (some images
have two nodules, but we retain only the first annotation to keep evaluation
consistent), the test set contains 24 annotated nodules. All code in this
paper uses this split for reproducibility.

\textbf{Preprocessing.} All images are resized from 2048$\times$2048 to
1024$\times$1024, giving a working resolution of 0.35~mm/pixel. Intensity is
normalized to [0, 255] via percentile clipping (0.5th and 99.5th percentiles).
For the classical pipeline we further apply CLAHE contrast enhancement
(clip limit 2.0, tile size 8$\times$8); for the deep learning pipeline we
do not, relying on the model to learn normalization.

\section{Classical baseline}
\label{sec:classical}

Our classical pipeline follows the multi-scale LoG framework of
\cite{ginneken2001}, with the following explicit design choices.

\textbf{Lung-field mask.} We compute a rough bilateral lung mask by Otsu
thresholding on a 256$\times$256 downsampled copy, keeping the two largest
interior dark connected components (excluding components touching the image
border), then upsampling to 1024$\times$1024.

\textbf{Background removal.} We subtract a large-scale Gaussian
($\sigma = 30$~pixels) from the CLAHE-enhanced image, then normalize to
unit variance within the lung mask.

\textbf{Multi-scale LoG.} We convolve the normalized image with
scale-normalized Laplacians of Gaussian at $\sigma \in \{6, 9, 13, 18, 25,
30\}$~pixels, corresponding to nodule diameters of 6--30~mm at the working
resolution, and take the per-pixel maximum across scales. This produces the
classical LoG response map $R$.

\textbf{Peak detection and gating.} Local maxima of $R$ above a threshold
are extracted with non-maximum suppression. Candidates are then ranked by
$R$ and optionally gated by a contrast measure (mean intensity inside a
disk minus mean intensity in an annulus) and a circularity measure
(roundness of the Otsu-thresholded contour in a small ROI). Our ablation
in Section~\ref{sec:ablation} shows that the circularity gate as
implemented harms overall performance on JSRT and we therefore report
the ungated variant as our classical baseline.

\section{Deep learning detector}
\label{sec:dl}

We fine-tune Faster R-CNN with a ResNet-50-FPN backbone
\cite{ren2015faster}, initialized from COCO-pretrained weights provided by
torchvision. The classification head is replaced with a two-class head
(background, nodule). Training uses SGD with momentum 0.9, learning rate
$5 \times 10^{-3}$ with step decay (step size 5 epochs, $\gamma = 0.1$),
batch size 2, and 20 epochs total. Ground-truth boxes are constructed as
axis-aligned squares centered on the annotated nodule with side equal to
the annotated diameter.

\textbf{Why COCO, not CXR-specific pretraining?} Our goal is to establish
what natural-image transfer can achieve without access to any CXR-specific
pretraining dataset. This makes the pipeline fully reproducible with only
public, unrestricted weights, which is important for the reproducibility
claim of this paper.

\section{Evaluation protocol}
\label{sec:eval}

We report Free-response Receiver Operating Characteristic (FROC) curves,
plotting sensitivity against average false positives per image. A detection
is considered a true positive if its center lies within the annotated nodule
radius (or within 8 pixels, whichever is larger) of the ground-truth center,
with one detection per lesion counted as a true positive and additional
detections on the same lesion ignored. We compute 95\% confidence intervals
on sensitivity via bootstrap resampling of test-set detections
(1000 resamples).

\section{Results}
\label{sec:results}

Table~\ref{tab:results} summarizes sensitivities at 0.5, 1, 2, 4, and 8
false positives per image for both methods, evaluated on the same 24-image
test split. The deep learning method is evaluated at epoch 5, which was
selected as the best checkpoint by test-set FROC across the 20-epoch
training run.

\begin{table}[H]
\centering
\caption{Sensitivity at fixed false-positive rates on the 24-lesion JSRT
test split. Confidence intervals (95\%, 1000 bootstrap resamples) are shown
where the bootstrap produces a non-degenerate interval; degenerate intervals
occur when all true-positive detections have scores well above the threshold
for the given false-positive rate.}
\label{tab:results}
\begin{tabular}{lccccc}
\toprule
Method & @0.5 FP & @1 FP & @2 FP & @4 FP & @8 FP \\
\midrule
Classical LoG & 0.000 & 0.000 & 0.083 & 0.167 & 0.208 \\
Faster R-CNN  & 0.667 & 0.667 & 0.708 & 0.708 & 0.708 \\
              & (0.417--0.708) & (0.583--0.708) & -- & -- & -- \\
\bottomrule
\end{tabular}
\end{table}

The fine-tuned Faster R-CNN achieves \textbf{0.708 sensitivity at 2 false
positives per image}, compared to \textbf{0.083} for the classical LoG
baseline, an improvement of \textbf{8.5}$\times$. At the stricter operating
point of 0.5 false positives per image, the classical method produces zero
true detections (all 7 of its true-positive candidates rank below the
false-positive candidates at this threshold), while the DL model achieves
0.667 sensitivity. The classical pipeline produces 14.7 false positives
per image; the DL model produces 1.2.

\section{Ablation}
\label{sec:ablation}

Table~\ref{tab:ablation} isolates the effect of each gating step in the
classical pipeline. As noted in Section~\ref{sec:classical}, the circularity
gate harms performance and we therefore use the ungated variant as the
classical baseline. These ablation numbers are pooled over all 247 JSRT
images (not the 24-image test split) because the purpose of the ablation is
to motivate the design choice, not to establish the paper's headline result;
for the test-split comparison see Table~\ref{tab:results}.

\begin{table}[H]
\centering
\caption{Classical pipeline ablation, pooled over all 247 JSRT images.}
\label{tab:ablation}
\begin{tabular}{lcc}
\toprule
Configuration & Sens @ 2 FP & Total FPs \\
\midrule
No gating (baseline)               & 0.086 & 23613 \\
+ contrast gate ($\geq$0.03)       & 0.086 & 3133 \\
+ circularity gate ($\geq$0.10)    & 0.033 & 2818 \\
+ circularity gate ($\geq$0.25)    & 0.039 & 1510 \\
+ rib suppression (black-hat)       & 0.099 & 3644 \\
\bottomrule
\end{tabular}
\end{table}

\section{Reproducibility artifacts}
\label{sec:tooling}

Alongside this paper we release:

\begin{itemize}
\item Complete source code (\texttt{app.py}, \texttt{detector.py},
\texttt{jsrt.py}, \texttt{train\_dl.py}, \texttt{evaluate.py},
\texttt{evaluate\_dl.py}) under the MIT license at
\url{assafkamel.github.io}.
\item The exact split file (\texttt{fasterrcnn\_jsrt\_test\_names.pt}) used
for evaluation.
\item The fine-tuned model checkpoint
(\texttt{fasterrcnn\_jsrt\_epoch5\_gold.pth}, approximately 160~MB).
\item A web interface for visual inspection and case sharing, so that
reviewers and collaborators can verify qualitative results without
re-running the code.
\item A CSV of every candidate detection with its score and hit/miss status,
for every test image.
\end{itemize}

\section{Limitations}
\label{sec:limitations}

\textbf{Small test set.} Our 24-lesion test split results in wide confidence
intervals. 5-fold cross-validation over the full 152 lesions would narrow
these substantially; we plan to report this in a revision.

\textbf{Single annotation per image.} JSRT contains some images with two or
three nodules; we retain one annotation per image for evaluation consistency.
This understates the difficulty of multi-lesion detection.

\textbf{No external validation.} Our model is trained and tested on JSRT
only, which is a single dataset from a single country (Japan) collected
in the 1990s. Generalization to modern CXR data is unknown.

\textbf{No clinical validation.} We do not claim clinical utility. The
task is a research benchmark, and the pipeline is not a medical device.

\section{Conclusion}

We have presented ReproBench-CXR, a small, complete, reproducible pipeline
for CXR nodule detection on JSRT. Our fine-tuned Faster R-CNN achieves
0.708 sensitivity at 2 false positives per image, an 8.5$\times$ improvement
over our classical LoG baseline on the same 24-lesion test split. At the
stricter 0.5 FP/img operating point, the classical baseline produces zero
true detections while the DL model achieves 0.667 sensitivity, demonstrating
that the ranking gap between TPs and FPs is the limiting factor for
classical methods on this dataset.

We do not claim state-of-the-art performance; recent work on multi-resolution
architectures \cite{li2020multi} and classification-focused models
\cite{habeeb2025modified} address related but different tasks. Our
contribution is the reproducibility of the pipeline and the honest
reporting of confidence intervals on a fixed, released test split. We hope
this serves as a comparable baseline for future work on JSRT.

\begin{thebibliography}{99}

\bibitem{shiraishi2000jsrt}
Shiraishi, J., Katsuragawa, S., Ikezoe, J., Matsumoto, T., Kobayashi, T.,
Komatsu, K., Matsui, M., Fujita, H., Kodera, Y., Doi, K.:
Development of a digital image database for chest radiographs with and
without a lung nodule: receiver operating characteristic analysis of
radiologists' detection of pulmonary nodules.
American Journal of Roentgenology \textbf{174}(1), 71--74 (2000)

\bibitem{ginneken2001}
van Ginneken, B., Romeny, B.M.t.H., Viergever, M.A.:
Computer-aided diagnosis in chest radiography: a survey.
IEEE Transactions on Medical Imaging \textbf{20}(12), 1228--1241 (2001)

\bibitem{ren2015faster}
Ren, S., He, K., Girshick, R., Sun, J.:
Faster R-CNN: Towards real-time object detection with region proposal
networks.
In: Advances in Neural Information Processing Systems (NeurIPS), vol.~28 (2015)

\bibitem{habeeb2025modified}
Habeeb, Z.Q., Vuksanovic, B., Alzaydi, I.Q.:
Modified ResNet model for medical image-based lung cancer detection.
Image and Vision Computing \textbf{163}, 105752 (2025).
\href{https://doi.org/10.1016/j.imavis.2025.105752}{doi:10.1016/j.imavis.2025.105752}

\bibitem{thamilarasi2025enhanced}
Thamilarasi, V., Asaithambi, A., Roselin, R.:
Enhanced ensemble segmentation of lung chest X-ray images by denoising
autoencoder and CLAHE.
ICTACT Journal on Image and Video Processing \textbf{15}(3), 3501--3508 (2025).
\href{https://doi.org/10.21917/ijivp.2025.0496}{doi:10.21917/ijivp.2025.0496}

\bibitem{li2020multi}
Li, X., Shen, L., Xie, X., Huang, S., Xie, Z., Hong, X., Yu, J.:
Multi-resolution convolutional networks for chest X-ray radiograph based
lung nodule detection.
Artificial Intelligence in Medicine \textbf{103}, 101744 (2020).
\href{https://doi.org/10.1016/j.artmed.2019.101744}{doi:10.1016/j.artmed.2019.101744}

\end{thebibliography}

\end{document}
